\scanpage{94}
\mspage{87}
% archive sheet number: 93
si les conditions suivantes sont satisfaites.

Pour tt \deleted{\unsure{$X\in\operatorname{Ob}E$, et}} système inductif $(X_i)_{i\in I}$ parfaitement adapté à la filtration cardinale pour $\pi$ (donc $I$ ens. ordonné \deleted{grand devant $\pi$;} \inserted{filtrant} tel que tt partie filtrante $\unsure{J|I}$ de cardinal $\leq\pi$ ait une borne supérieure $i_J$, et $X_{i_J}\xleftarrow{\simeq}\varinjlim_{j\in J}X_j$), \deleted{de limite inductive $X$ \uline{ds $M$},} \deleted{\unsure{alors pour tt $i\in I$ on ait}.} \inserted{Soit $I_M$ l'ens des} \deleted{\unsure{l'ens}} \deleted{$I_M$ des} $i\in I$ tels que $X_i\in M$. \deleted{est cofinal dans $I$} Ceci posé, on veut ceci:

a) pour que $X$ ($=\varinjlim X_i$) soit dans $M$, il \unsure{faut (et) suffit|faut et suffit}, \deleted{\unsure{par définition}} \deleted{\unsure{que}} \inserted{que} $I_M$ soit cofinal dans $I$ (de sorte que $X$ apparaisse comme $\varinjlim$ filtrante, d'un ens. d'indices $I_M$ grand devant $\pi$, d'éléments de $M\cap\mathrm{Filt}^{\pi}\overset{\mathrm{def}}{=}M^{\pi}$).

\deleted{b) $M$ est stable par \deleted{petites} limites \unsure{inductives} filtrantes \inserted{(\unsure{en $\underline M$}),} grandes devant $\pi$.
% slanted
\margin{\unsure{Découle|Donc} \uline{donc} la condition b), qui est indépendante d'un syst. inductif particulier.}

On dit que $M$ est accessible s'il existe un cardinal $\pi\geq\pi_0$ tel que $M$ soit $\pi$-accessible.}
